\begin{table*}[t]
\centering
\small
\setlength{\tabcolsep}{4pt}
\begin{tabularx}{\linewidth}{>{\raggedright\arraybackslash}p{2.5cm}>{\raggedright\arraybackslash}p{2.9cm}>{\raggedright\arraybackslash}p{2.7cm}>{\raggedright\arraybackslash}X}
\toprule
Benchmark & Input and task & Evaluation signal & Relative limitation \\
\midrule
FinQA / ConvFinQA / TAT-QA / DocFinQA
& No raw transactions; question answering over reports, tables, or filings
& Partial or answer-level scoring
& Does not evaluate statement construction, global accounting invariants, or closing behavior \\
FinanceMATH / MMLU
& No raw transactions; financial reasoning or exam-style QA
& Exact answer scoring
& Tests quantitative reasoning, but not bookkeeping structure or executable balance-sheet assembly \\
FinanceBench / FinBen
& No raw transactions; analysis over prepared financial artifacts
& Mixed evaluation settings
& Focuses on analysis and judgment rather than constructing a statement under hard constraints \\
FinAuditing / FinRule-Bench / auditing benchmarks
& Mixed inputs; auditing, diagnosis, or rule-checking tasks
& Exact or partially structural scoring
& Closer to accounting workflows, but still does not require full balance-sheet construction with closing evaluation \\
\benchmark{} (ours)
& Raw journal events transformed into a complete balance sheet
& Exact account-level scoring plus invariant and closing checks
& Directly evaluates executable statement construction from transactions \\
\bottomrule
\end{tabularx}
\caption{Positioning relative to prior work. Existing financial benchmarks mostly evaluate question answering, analysis, or auditing over prepared artifacts. \benchmark{} instead starts from raw journal events and evaluates executable statement construction under exact accounting invariants, including retained-earnings closure.}
\label{tab:prior-work-comparison}
\end{table*}
